\chapter{Partial Differential Equations}

% -------------------------------------------------------
\section{General Second-Order Form}
% -------------------------------------------------------

The general second-order linear PDE in two independent variables $x$ and $y$ is
\[
  a\,u_{xx} + b\,u_{xy} + c\,u_{yy} + d\,u_x + e\,u_y + f\,u = g(x,y),
\]
where $a$, $b$, $c$, $d$, $e$, $f$ are (possibly variable) coefficients and subscripts denote partial derivatives.

% -------------------------------------------------------
\section{Classification}
% -------------------------------------------------------

\begin{theorem}[Classification by Discriminant]
The type of the second-order PDE is determined by the \textbf{discriminant} $\Delta = b^2 - 4ac$:
\end{theorem}

\begin{table}[h!]
\centering
\renewcommand{\arraystretch}{1.4}
\begin{tabular}{lcc}
\toprule
\textbf{PDE Type} & \textbf{Condition} & \textbf{Physical Analogue} \\
\midrule
Elliptic   & $b^2 - 4ac < 0$ & Steady-state (e.g.\ Laplace) \\
Parabolic  & $b^2 - 4ac = 0$ & Diffusion (e.g.\ Heat equation) \\
Hyperbolic & $b^2 - 4ac > 0$ & Wave propagation \\
\bottomrule
\end{tabular}
\caption{Classification of second-order PDEs.}
\end{table}

% -------------------------------------------------------
\section{Canonical Examples}
% -------------------------------------------------------

\subsection{Elliptic: Laplace's Equation}

\[
  u_{xx} + u_{yy} = 0
\]

Here $a=1$, $b=0$, $c=1$, so $b^2 - 4ac = 0 - 4 = -4 < 0$. \textbf{Elliptic.}

\begin{remark}
Laplace's equation governs steady-state temperature distributions, electrostatics, and irrotational fluid flow.
\end{remark}

\subsection{Parabolic: Heat Equation}

\[
  u_t = \beta\, u_{xx}
\]

Rewriting as $-\beta\,u_{xx} + u_t = 0$, we identify $a = \beta$, $b = 0$, $c = 0$ (treating $t$ as the second spatial variable), giving $b^2 - 4ac = 0$. \textbf{Parabolic.}

\begin{remark}
The heat equation models diffusion processes: heat conduction, mass diffusion, and probability densities in stochastic processes.
\end{remark}

\subsection{Hyperbolic: Wave Equation}

\[
  u_{tt} = \alpha^2\, u_{xx}
\]

With $a = -\alpha^2$, $b = 0$, $c = 1$ (identifying independent variables $x$ and $t$), we get $b^2 - 4ac = 0 + 4\alpha^2 > 0$. \textbf{Hyperbolic.}

\begin{remark}
The wave equation describes vibrating strings, acoustic waves, and electromagnetic wave propagation.
\end{remark}
